\documentclass[conference,twocolumn,a4paper]{IEEEtran}
\usepackage{amssymb,amsmath,amsthm} % math packages 
\interdisplaylinepenalty=2500
\usepackage{cite} % citation 
\usepackage{psfrag}
\usepackage{here}
\usepackage{algorithm}
\usepackage{algpseudocode}
\algblock{Output}{EndOutput}
\algnotext{EndOutput}
\algnewcommand{\Initialize}[1]{%
	\State \textbf{Initialize:}
	\Statex \hspace*{\algorithmicindent}\parbox[t]{.8\linewidth}{\raggedright #1}
}
%\usepackage{graphicx,subfigure} % figures 
\usepackage[left=1.57cm, right=1.58cm, top=1.71cm,bottom=4.30cm]{geometry}

\usepackage{graphicx,psfrag}
\usepackage{fancyhdr,afterpage,lipsum}
\usepackage{booktabs}
\usepackage{tabularx}
\usepackage{multicol}
\usepackage{multirow}
\usepackage{color}
\usepackage{subcaption}
\usepackage{url} % links 
\usepackage[latin1]{inputenc} % language
\DeclareMathOperator*{\argmmin}{arg\,min}
\DeclareMathOperator*{\argmmax}{arg\,max}

\newtheorem{proposition}{Proposition} 
\theoremstyle{remark} \newtheorem{remark}{Remark} 

\title{Multi-Static ISAC in Cell-Free Massive MIMO: Precoder Design and Privacy Assessment}

\author{\IEEEauthorblockN{Isabella W. G. da Silva, Diana P. M. Osorio, and Markku Juntti}
\IEEEauthorblockA{\textit{Centre for Wireless Communications, University of Oulu, P.O.Box 4500, FI-90014, Finland} \\
Emails: \{isabella.dasilva, diana.moyaosorio, markku.juntti\}@oulu.fi}}


\begin{document}
\maketitle
\vspace{-1cm}
\begin{abstract}
%\boldmath
A multi-static sensing-centric integrated sensing and communications (ISAC) network can take advantage of the cell-free massive multiple-input multiple-output infrastructure to achieve remarkable diversity gains and reduced power consumption. While the conciliation of sensing and communications requirements is still a challenge, the privacy of the sensing information is a growing concern that should be seriously taken on the design of these systems to prevent other attacks. This paper tackles this issue by assessing the probability of an internal adversary to infer the target location information from the received signal by considering the design of transmit precoders that jointly optimizes the sensing and communications requirements in a multi-static-based cell-free ISAC network. Our results show that the multi-static setting facilitates a more precise estimation of the location of the target than the mono-static implementation. 
\end{abstract}
\begin{IEEEkeywords}
cell-free massive MIMO, ISAC, multi-static sensing, precoder design, privacy. 
\end{IEEEkeywords}
\section{Introduction}
The sixth generation (6G) of wireless communications is expected to heavily influence every aspect of the Society of 2030. Several applications envisioned for 6G will be supported by sensing capabilities, with the network acting as a sensor. For this becoming reality, the new paradigm of integrated sensing and communications (ISAC) plays an important role by allowing  intelligent sharing of wireless resources and the support for a number of use cases such as autonomous vehicles and smart homes~\cite{9737357}.  

So far, the research on ISAC systems has mainly focused on mono-static sensing. However, the employment of mono-static sensing requires the co-located sensing transmitter and receiver to be full-duplex. Hence, multi-static sensing ISAC systems, in which there are multiple non-colocated transmitters and receivers for sensing, can provide diversity gain while avoiding the need for full-duplex nodes~\cite{Art:Emil}. For instance, in~\cite{Art:Emil}, Behdad \textit{et al.} assessed the probability of detection of a single target in a cell-free massive multiple-input multiple-output MIMO ISAC system where a number of ISAC transmitters and sensing receivers are deployed to detect the target and to transmit data to multiple communications users (UEs). Moreover, in~\cite{9842350}, Huang \textit{et al.} proposed a coordinated power control design for a networked ISAC system aiming to maximize the signal-to-interference-plus-noise ratio (SINR) at the UEs and maximize the Cr\'{a}mer-Rao lower bound (CRLB) of the target location estimate. Both works demonstrated that, in comparison to communication-centric designs, a cell-free massive MIMO ISAC design requires less transmit power to attain more accurate detection.

Nevertheless, the design of joint precoders for sensing and communications employed in most of works with mono-static or multi-static ISAC systems considered the simultaneous transmission of individual radar and communications waveforms~\cite{9124713,demirhan2023cell}. Thus, the transmit precoder contains information about the target location and the data intended for the UEs. Indeed, the advancements on ISAC design have exacerbated security and privacy concerns since the communications information could be exposed to untrusted targets, and malicious UEs could try to infer the position of targets. In this sense, security and privacy aspects have recently started to emerge as a crucial part for the design of ISAC. For instance, the works in \cite{art:su2021, 10104579} and \cite{art:ren} have focused on the design of beamformers with secrecy requirements to prevent untrusted targets from eavesdropping information from UEs. In~\cite{art:su2021}, Su \textit{et al.} proposed a beamforming design to minimize the signal-to-noise ratio (SNR) at a potential malicious target constrained to minimal SINR requirements for the UEs. In~\cite{10104579}, the transmit beamforming is obtained via a weighted optimization between radar estimation CRLB and the communications secrecy rate. Also, in~\cite{art:ren}, Ren \textit{et al.} proposed a beamforming design to minimize the beampattern matching error constrained to secrecy rates requirements in a network with multiple targets, with some of them assumed as malicious.

On the other hand, the privacy of the sensing information remains barely explored while being an important source of concern as it can trigger further attacks. In our earlier work~\cite{dasilva2023privacy}, we addressed this aspect by evaluating the capability that an internal adversary of the network will have to infer the angular position of a target by an UE that acts as an adversary. Results demonstrated that if the adversary is capable to infer the transmit precoder, the detection surpasses fifty percent of the cases, which entails a serious privacy breach for ISAC systems. 

Accordingly, recognizing the benefits of multi-static sensing and the seriousness of possible privacy leakages in ISAC systems, in this work we investigate a multi-static sensing-centric ISAC network in a cell-free massive MIMO scenario, with multiple MIMO UEs and a point-like target. In this scenario, one of the UEs acts maliciously by trying to estimate the position of the target from the observations of the received signal. The main contributions of this work are as follows: i) we propose a transmit precoder design to maximize the sensing SINR constrained to a minimal SINR for the UEs and a transmit power limit requirement; ii) a replica of the transmit beampattern from each transmitter access point (AP) is computed based on an expectation-maximization (EM) approach. 
%schemes to prevent the eavesdropping from untrusted targets have been designed in~\cite{art:su2021} and~\cite{art:ren}. Su \textit{et al.} proposed a beamforming design in~\cite{art:su2021} to minimise the signal-to-noise ratio (SNR) at the eavesdropper target while ensuring a minimal signal-to-interference-plus-noise ratio (SINR) at each CU and beampattern requirements. Moreover, Ren \textit{et al.} considered a network with multiple sensing targets in~\cite{art:ren}, where a portion of them are untrusted, thus a beamforming design was proposed to minimise the beampattern matching error restricted to secrecy rates requirements.   

%\textit{Notation:} Throughout this paper, bold upper-case letters denote matrices whereas bold lower-case letters denote vectors; $(\cdot)^T$, $(\cdot)^H$, and $(\cdot)^*$ stands for the matrix transpose, the Hermitian transpose, and the complex conjugate, respectively; $\mathbf{I}$ is the identity matrix; $||\cdot||$ and $|\cdot|$ are the Euclidean-norm and the absolute value operator; $\mathrm{Tr}(\cdot)$ is the trace of a square matrix; $p_X$ is the probability density function (PDF) of the variable $X$, $p_{X,Y}(\cdot)$ and $p_{X|Y}(\cdot)$, in particular, stands for the joint PDF of $X$ and $Y$, and the conditional PDF of $X$ and $Y$, respectively; $\mathbb{E}[\cdot]$ is the expectation operator; and $\mathcal{G}(\cdot,\cdot)$ is the gamma function.
\section{System Model}\label{sec:model}
\vspace{-0.25cm}
\begin{figure}[h]
    \centering
    \includegraphics[scale=0.2]{Figures/figmodel(1)(1).png}
    \caption{System Model}
    \label{fig:systemmodel}
  \vspace{-0.25cm}
\end{figure}

Consider the ISAC system as illustrated in Fig.~\ref{fig:systemmodel}, where $N_{\mathrm{Tx}}$ APs are jointly serving $N_{\mathrm{Ue}}$ UEs equipped with $M_{\mathrm{Ue}}$ antennas, while also detecting a point-like target. In this system, one of the UEs acts as an adversary and intends to discover the location of the target. The other $N_{\mathrm{Rx}}$ APs are acting as sensing receivers that simultaneously sense the location of the target. The transmit and receive APs are equipped with $M_{\mathrm{AP}}$ half-wavelength-spaced isotropic antennas deployed as an uniform linear array (ULA). It is further assumed that all APs are connected to the edge cloud, where the processing is done in a centralized manner, and are fully synchronized. Similar to~\cite{Art:Emil}, we assume that the transmitted signal by the $k$th transmit AP is given by a weighted sum of communications symbols and sensing signals. Hence, at time instance $m$, the transmitted signal $\mathbf{x}_k[n]$$\in$$\mathbb{C}^{M_{\mathrm{AP}}}$ can be written as  
\vspace{-0.2cm}
\begin{align}
    \mathbf{x}_k[n]\!=\!\sum_{i=1}^{N_{\mathrm{Ue}}} \mathbf{w}_{i,k} s_i[n]\!+\!\mathbf{w}_{t,k} s_t[n]\!=\!\mathbf{W}_k \mathbf{s}[n],\nonumber\\ n = 0,..., N\!-\!1, k = 1,..., N_{\mathrm{Tx}},
\end{align}
where $\mathbf{s}[n]$$=$$[s_1[n],...,s_{N_{\mathrm{Ue}}},s_{t}[n]]^T$ is the $(N_{\mathrm{Ue}}\!+\!1)$$\times$$1$ vector containing the $N_{\mathrm{Ue}}$$\times$$1$ parallel communications symbols intended to the $N_{\mathrm{Ue}}$ users plus the sensing signal, which is independent of the UE's data signals. Also, $\mathbf{w}_{i,k}$ and $\mathbf{w}_{t,k}$ are the $M_{\mathrm{AP}}\!\times\! 1$  transmit precoder vector of the $k$th transmitter AP for the $i$th user and for the sensing of the target, respectively. In addition, the users are assumed to employ a receive beamformer $\mathbf{u}_{i}$, with size $M_{\mathrm{Ue}}$$\times$$1$, to estimate the transmitted data stream. Accordingly, the estimated data stream at the $i$th user at time $n$ is given by
\begin{align}\label{eq:signali}
    y_{i}[n]&= \mathbf{u}_i^H\left(\sum_{k=1}^{N_{\mathrm{Tx}}}\mathbf{H}_{i,k}\mathbf{W}_{k}\mathbf{s}[n] +\mathbf{n}_i[n]\right),
\end{align}
where $\mathbf{H}_{i,k}$ is the $M_{\mathrm{Ue}}$$\times$$M_{\mathrm{AP}}$ channel coefficient matrix between the $k$th AP and UE $i$, which undergoes Rayleigh block fading. Moreover, $\mathbf{n}_i$ is the noise component at the $i$th user, modeled as signal-independent, zero-mean, additive white Gaussian noise (AWGN) with variance $\sigma^2_i\mathbf{I}$. Following, the SINR at UE $i$ given the transmission of AP $k$ is written as
\begin{align}
    \gamma_{i,k}\!=\!\frac{|\mathbf{u}_i^H\mathbf{H}_{i,k}\mathbf{w}_{i,k}|^2}{\sum_{l=1 \atop l \neq i}^{N_{\mathrm{Ue}}}\!|\mathbf{u}_i^H\mathbf{H}_{i,k}\mathbf{w}_{l,k}|^2\!+\!|\mathbf{u}_i^H\mathbf{H}_{i,k}\mathbf{w}_{t,k}|^2\!+\!\sigma^2_i||\mathbf{u}_i||^2}.
\end{align}
On the other hand, the signal received by the $r$th receiver AP at the time instant $n$ is written as
%Furthermore, since the sensing is deployed in a multi-static manner, the $N_{\mathrm{Rx}}$ APs serving as sensing receivers are not co-located with the transmitter APs. Each receiver AP receives desired reflected signals from the target when it is present, but also receives clutter signals composed of non-line-of-sight (NLOS) paths due to permanent and temporary obstacles. Permanent obstacles includes buildings and mountains, and the paths caused by these obstacles can be estimated beforehand. Thus, we assume that that the clutter caused by known obstacles is known at the edge cloud and can be canceled from the received signal at the receiver APs. On the other hand, the NLOS paths caused by temporary obstacles remain unknown, but can be neglected  if the target has good reflectivity and the interference originated by these unknown paths is assumed low~\cite{Art:Emil}. Accordingly, under these considerations and given that the target is present, the signal received by the $r$th receiver AP at the time instant $n$ is written as
\begin{align}
\mathbf{y}_r[n]\!=\!\!\!\sum_{k=1}^{N_{\mathrm{Tx}}}\!\!\alpha_{r,k} \sqrt{\beta_{r,k}}\mathbf{a}(\phi_r)\mathbf{a}^T(\phi_k)\mathbf{x}_k[n] \!+\!\mathbf{n}_r[n],
\end{align}
where $\alpha_{r,k}$ is the bi-static unknown radar cross section (RCS) of the target through the reflection path from transmitter AP $k$ to the receiver AP $r$. Herein, as in~\cite{Art:Emil}, it is assumed that the RCS follows the Swerling-I model, so that $\alpha_{r,k}$ is constant throughout a collection of consecutive sensing signals and follows the distribution $\alpha_{r,k}$$\sim$$\mathcal{CN}(0,\sigma_{r,k}^2)$. Also, $\phi_i$ is the azimuth angle from the target position to the $i$th AP, with $i \in \{k = 1, \dots, N_{\mathrm{Tx}},r = 1, \dots, N_{\mathrm{Rx}}\}$. 
%respectively. Whereas, $\phi_k$ and $\theta_k$ are the azimuth and elevation angles from the $k$th transmitter AP to the target location, respectively. 
Accordingly, $\mathbf{a}(\phi) \in \mathbb{C}^{M_{\mathrm{AP}}}$ is the antenna array steering vector, given by
\begin{align}
    \mathbf{a}(\phi)\!\!=\!\!\left[
        1, e^{j\pi\sin(\phi)}, \dots, e^{j(M_{\mathrm{AP}}\!-\!1)\pi\sin(\phi)}\right]^T\!\!\!\!.
\end{align}
Finally, $\mathbf{n}_r[n]$$\sim$$\mathcal{CN}(\mathbf{0},\sigma_n^2\mathbf{I}_{M_{\mathrm{AP}}})$ is the noise component at the $r$th receiver AP, and $\beta_{r,k}$ is the channel gain of the path between the $k$th transmitter AP to the target and from the target to the $r$th receiver AP, which can be calculated as~\cite{Art:Emil} %~\cite[Chapter~2]{book:radar}
\begin{align}
    \beta_{r,k} = \frac{\lambda_c^2}{(4\pi)^3d_{t,k}^2d_{t,r}^2},
\end{align}
where $\lambda_c$ is the carrier wavelength, and $d_{t,i}$ is the distance between the target location and the AP i, with $i \in \{k = 1, \dots, N_{\mathrm{Tx}},r = 1, \dots, N_{\mathrm{Rx}}\}$. Moreover, the possible clutter caused by permanent or temporary objects is neglected. Accordingly, the sensing SINR, $\gamma_t$, is given by%\footnote{For more details, please refer to~\cite{Art:Emil}.}
\begin{align}\label{eq:gammat}
    \gamma_t\!=&\zeta\!\Bigg(\!\sum_{n=1}^{N-1}\!\!\mathbf{s}^H\![n]\!\Bigg(\!\sum_{r=1}^{N_{\mathrm{Rx}}}\! \sum_{k=1}^{N_{\mathrm{Tx}}}\! \sum_{j=1}^{N_{\mathrm{Tx}}}\! \varphi_{r,k,j}\mathbf{W}_k^H\mathbf{A}_{k,j}\mathbf{W}_j\!\Bigg) \!\mathbf{s}[n]\!\Bigg)\!,
\end{align}
with $\varphi_{r,k,j} = \sqrt{\beta_{r,k} \beta_{r,j}}\mathbf{a}^H\left(\phi_r\right)\operatorname{cov}\left(\alpha_{r, j}, \alpha_{r, k}\right) \mathbf{a}\left(\phi_r\right)$, $\mathbf{A}_{k,j}$$=$$\mathbf{a}^*\!\left(\phi_k\right)\mathbf{a}^T\!\left(\phi_j\right)$, and $\zeta$$=$$1/(N\!-1\!)M_{\mathrm{AP}}N_{\mathrm{Rx}}\sigma_n^2$.
\section{Precoder Design}
For the design of the ISAC transmit precoding, the goal is to maximize the sensing SINR in \eqref{eq:gammat} under transmit power and communications quality of service (QoS) constraints, which can be formulated as
\begin{subequations}
\begin{align}
\mathcal{P}: \max_{\mathbf{W},\mathbf{u}} &  \hspace{3mm}\gamma_t \label{eq:objP}\\
\text { s. t. } & \gamma_{i,k} \geq \Gamma, i=1, \ldots, N_{\mathrm{Ue}}, k=1, \ldots, N_{\mathrm{Tx}}\label{eq:gammai}\\
& \sum_{i=1}^{N_{\mathrm{Ue}}}||\mathbf{w}_{i,k}||^2\!+\!||\mathbf{w}_{t,k}||^2\!\leq \!P_T, k=1,...,N_{\mathrm{Tx}},\label{eq:pt}
\end{align}
\end{subequations}
where $\Gamma$ is the communications SINR threshold, and $P_T$ is the transmit power limit for the transmitter APs. To solve $\mathcal{P}$, an iterative process is employed. At first, the UEs receive beamformers $\mathbf{u}_i$, $i$$=$$1,...,N_{\mathrm{Ue}}$ are randomly chosen and fixed. Nonetheless, note that the optimization problem in $\mathcal{P}$ is still not convex due to the non-concave objective function and non-convex constraint in \eqref{eq:gammai}. To handle this, \eqref{eq:objP} can be split into two sub-problems as
\begin{subequations}
\begin{align}
\mathcal{P}1\!:\! \max_{\mathbf{W}} &\zeta\!\Bigg(\!\sum_{n=1}^{N\!-\!1}\!\!\mathbf{s}^H\![n]\!\Bigg(\!\sum_{r=1}^{N_{\mathrm{Rx}}}\! \sum_{k=1}^{N_{\mathrm{Tx}}}\! \varphi_{r,k,k}\mathbf{W}_k^H\mathbf{A}_{k,k}\mathbf{W}_k\!\Bigg)\!\mathbf{s}[n]\!\!\Bigg)\label{eq:p1}\\
\text { s. t. } & \eqref{eq:gammai}, \eqref{eq:pt},\nonumber 
\end{align}
\end{subequations}
and, 
\begin{subequations}
\begin{align}
\mathcal{P}2\!: \max_{\mathbf{W}} &\zeta\!\Bigg(\!\sum_{n=1}^{N-1}\!\!\mathbf{s}^H\![n]\!\Bigg(\!\sum_{r=1}^{N_{\mathrm{Rx}}}\! \sum_{k=1}^{N_{\mathrm{Tx}}}\! \sum_{j=1 \atop j \neq k}^{N_{\mathrm{Tx}}}\! \varphi_{r,k,j}\mathbf{W}_k^H\mathbf{A}_{k,j}\mathbf{W}_j\!\Bigg) \mathbf{s}[n]\!\Bigg)\label{eq:p2}\\
\text { s. t. } & \eqref{eq:gammai},\eqref{eq:pt}.\nonumber 
\end{align}
\end{subequations}
Next, as in~\cite{art:paperRFVLC}, we consider the first-order Taylor approximation to linearize \eqref{eq:p1} and \eqref{eq:gammai}, and employ the constrained convex-concave procedure (CCCP) to iteratively solve the approximated convex problem until a convergence criterion is attained. Thus, at the $p$th iteration, $\mathcal{P}1$ is solved as
%Next, as in~\cite{art:paperRFVLC}, we consider the first-order Taylor approximation to linearize the objective function of $\mathcal{P}1$, \eqref{eq:p1} and the constraint in \eqref{eq:gammai}, and employ the CCCP to iteratively solve the approximated convex problem until a convergence criterion is attained. Thus, at the $p$th iteration, $\mathcal{P}1$ is solved as
\begin{subequations}
\begin{align}
\mathcal{P}1'&\!:\! \max_{\mathbf{W}, \tau_{i,k}, \mu_{i,k}}\!\! \zeta\!\Bigg(\!\sum_{n=1}^{N\!-\!1}\!\!\mathbf{s}^H\![n]\!\Bigg(\!\sum_{r=1}^{N_{\mathrm{Rx}}}\! \sum_{k=1}^{N_{\mathrm{Tx}}}\! \varphi_{r,k,k}\!\Bigg(\!\!2\left[\!\mathbf{W}_k^{(p\!-\!1)}\!\right]^H\!\!\!\mathbf{A}_{k,k}\nonumber\\
\!\times&\!\!\Bigg(\!\mathbf{W}_k\!-\!\mathbf{W}_k^{(p\!-\!1)}\!\!\Bigg)\!\!+\!\!
\left[\!\mathbf{W}_k^{(p\!-\!1)}\!\right]^H\!\!\!\mathbf{A}_{k,k}\mathbf{W}_k^{(p\!-\!1)}\!\Bigg)\!\!\!\Bigg)\!\mathbf{s}[n]\!\Bigg)\label{eq:p1linha}\\
\text { s. t. }\!&|\mathbf{u}_i^H\mathbf{H}_{i,k}\mathbf{w}_{i,k}^{(p-1)}|^2\!+\!2\!\left[\!\mathbf{w}_{i,k}^{(p-1)}\!\right]^H\!\!\mathbf{H}_{i,k}^H\mathbf{u}_i\mathbf{u}_i^H\mathbf{H}_{i,k}\nonumber\\
&\times\!\left(\!\mathbf{w}_{i,k}\!-\!\mathbf{w}_{i,k}^{(p-1)}\!\right)\!\geq \tau_{i,k}, \forall i, \forall k\label{eq:SINR1}\\
&\sum_{l=1 \atop l \neq i}^{N_{\mathrm{Ue}}}\!|\mathbf{u}_i^H\mathbf{H}_{i,k}\mathbf{w}_{l,k}|^2\!+\!|\mathbf{u}_i^H\mathbf{H}_{i,k}\mathbf{w}_{t,k}|^2\!+\!\sigma^2_i||\mathbf{u}_i||^2\!\leq\!\mu_{i,k}\label{eq:SINR2}\\
&\tau_{i,k} \geq \Gamma \mu_{i,k}, \forall i, \forall k,\label{eq:slack}\\
&\eqref{eq:pt},\nonumber 
\end{align}
\end{subequations}
where $\tau$ and $\mu$ are slack variables. $\mathcal{P}1'$ is a concave problem. 
%that can be efficiently solved with convex toolboxes such as CVX.  
$\mathcal{P}2$, on the other hand, is approximated to the following
\begin{subequations}
\begin{align}
\mathcal{P}2': \max_{\mathbf{W}, \tau_{i,k},\mu_{i,k}} &\zeta\!\Bigg(\!\sum_{n=1}^{N-1}\!\!\mathbf{s}^H\![n]\!\Bigg(\!\sum_{r=1}^{N_{\mathrm{Rx}}}\! \sum_{k=1}^{N_{\mathrm{Tx}}}\! \sum_{j=1 \atop j \neq k}^{N_{\mathrm{Tx}}}\! \varphi_{r,k,j}\mathbf{W}_k^H\mathbf{A}_{k,j}\nonumber\\&\times\mathbf{W}_j^{(p-1)}\Bigg)\mathbf{s}[n]\Bigg)\label{eq:p21}\\
\text { s. t. } & \eqref{eq:SINR1}, \eqref{eq:SINR2},\eqref{eq:slack}, \eqref{eq:pt}.\nonumber 
\end{align}
\end{subequations}
which is a concave problem. Accordingly, both optimization problems can be solved with convex solvers such as CVX. After obtaining the optimal $p$th $\mathbf{W}^*$, given by the summation of the solutions of $\mathcal{P}1'$ and $\mathcal{P}2'$, the transmit precoder matrices are fixed and, for all $i$, $\mathbf{u}_i$ is updated via the MMSE receiver as~\cite{7405344}
\begin{align}
         \mathbf{u}_i\!&=\!\!\sum_{k=1}^{N_{\mathrm{Tx}}}\!\!\left(\!\mathbf{H}_{i,k}\!\!\left(\!\sum_{l=1 \atop l \neq i}^{N_{\mathrm{Ue}}}\mathbf{w}_{l,k}\mathbf{w}_{l,k}\!\!\right)\!\mathbf{H}_{i,k}^H\!+\!\sigma_i^2\mathbf{I}_{M_{\mathrm{Ue}}}\!\!\right)^{\!-\!1}\!\!\!\mathbf{H}_{i,k} \mathbf{w}_{i,k}. \label{eq:ui}
     \end{align}
     
Accordingly, the algorithm to obtain the transmit and receive precoders is summarized in Algorithm 1. 
\begin{algorithm}\label{alg:CCCP}
\caption{Precoders Iterative Algorithm}
\begin{algorithmic}[1]
\footnotesize
       \State Choose error tolerances $\epsilon$ , maximum number of iterations $p_{max}$, and feasible initial points $\mathbf{W}^{(0)}$, and $\mathbf{u}_i(0)$ $\forall i$.
       \State $p \gets 0$
       \Repeat
       \State Solve \eqref{eq:p1linha} and \eqref{eq:p21} to obtain $\mathbf{W}^{*(p)}$ using $\mathbf{W}^{*(p-1)}$ and $\mathbf{u}_{i}^{(p-1)}$ from the previous iteration.
       \State Solve \eqref{eq:ui} using $\mathbf{W}^{*(p)}$ to attain $\mathbf{u}_{i}^{(p)}$.
       \State $p++$
       \Until{$\frac{||\mathbf{W}^{*(p)}-\mathbf{W}^{*(p-1)}||}{\mathbf{W}^{*(p)}}$$\leq$$\epsilon$ and $\frac{||\mathbf{u}_i^{*(p)}-\mathbf{u}_i^{*(p-1)}||}{\mathbf{u}_i^{*(p)}}$$\leq$$\epsilon$ or $p$$\geq$$p_{max}$}
      \State \Return $\mathbf{W}^*$, $\mathbf{u}_i^*$ $\forall i$.  
\end{algorithmic}
\end{algorithm}
\section{Adversary Model}
Assuming that the UEs know that there is a target present in the network, one of them may act as an adversary. Given its received signal, the UE could try to infer the position of the target by recreating a replica of the transmit beampattern created by the transmitter APs given by
\begin{align}
    \mathbf{B}_k = \mathbf{a}^H(\theta_{z})\mathbf{R}_{\mathbf{x}_k}\mathbf{a}(\theta_{z}), \forall k,
\end{align}
where $\{\theta_z\}^{Z}_{z=1}$ are sampled angle grids, and $\mathbf{R}_{\mathbf{x}_k}$$=$$\mathbf{x}_k\mathbf{x}_k^H$ is the covariance matrix of the transmitted signal. To accomplish this, it is considered that the adversary has knowledge of its receive beamforming vector. Also, it is feasible to assume that the adversary knows the position of the $N_{\mathrm{Tx}}$ transmitter APs and can create a search area based on that information. Under these considerations, the adversary estimation process starts with the inference of the interference signal from its received signal, which is expressed as in \eqref{eq:signali}, is equal to
\begin{align}
     \tilde{\mathbf{y}}_{a,k} = \mathbf{H}_{a,k}\underbrace{\left(\sum_{l=1 \atop l \neq a}^{N_{\mathrm{Ue}}}\!\mathbf{w}_{l,k}s_l[n]+\mathbf{w}_{t,k}s_t[n]\right)}_{\tilde{\mathbf{x}}_k}+\mathbf{n}_a.
 \end{align}
We further assume that the true $\mathbf{H}_{a,k}$ is not available at the adversary, but can be estimated as $\mathbf{H}_{a,k}\!=\!\Hat{\mathbf{H}}_{a,k}\!+\!\varepsilon$, where $\Hat{\mathbf{H}}_{a,k}$ is its channel estimate, which can be attained via least squares estimation with pilot signals as 
\begin{align}
    \Hat{\mathbf{H}}_{a,k} = \mathbf{y}_{a,k}\mathbf{x}_k^H(\mathbf{x}_k\mathbf{x}_k^H)^{-1},
\end{align}
and $\varepsilon$ represents the channel estimation error, modeled as $\mathcal{N}(0,\sigma_\varepsilon^2 \mathbf{I})$. Thus, $\mathbf{H}_{a,k}$$\sim$$\mathcal{N}(\Hat{\mathbf{H}}_{a,k},\sigma_\varepsilon^2 \mathbf{I})$.

Then, considering  $\tilde{\mathbf{y}}_{a,k}$ as an observable variable, $\mathbf{H}_{a,k}$ as a latent variable, and   $\tilde{\mathbf{x}}_k$ as the unknown parameter, an iterative method to compute the maximum log-likelihood as the EM can be employed to estimate $\tilde{\mathbf{x}}_k$. The EM algorithm is a general technique used to find the maximum likelihood estimates of parameters when part of the data is missing or for latent variable models. The algorithm is divided into two steps, the E-step and the M-step. In the former, the expected value of the log-likelihood function is computed, as well as a current estimate of the unknown parameter. In the latter, the unknown parameter is attained from the maximization of the previously obtained expected value~\cite{art:em}.
\subsection{Expectation-Maximization Algorithm}
As the goal is to maximize the log-likelihood function, we start by defining it as
\begin{align}\label{eq:likelihood}
    \mathcal{L}\!=\!\log \left(p_{\tilde{\mathbf{y}}_{a,k}}\!({\tilde{\mathbf{x}}_k})\right)\!=\!\log\!\!\left(\int\!\! p_{\tilde{\mathbf{y}}_{a,k},\mathbf{H}_{a,k}}\!(\tilde{\mathbf{x}}_{k}) d\mathbf{H}_{a,k}\right), 
\end{align}
where $p_{\tilde{\mathbf{y}}_{a,k}}$ is the probability density function (PDF) of the variable $\tilde{\mathbf{y}}_{a,k}$, and $p_{\tilde{\mathbf{y}}_{a,k},\mathbf{H}_{a,k}}\!(\tilde{\mathbf{x}}_{k})$ stands for the joint PDF of $\mathbf{y}_{a,k}$ and $\mathbf{H}_{a,k}$, and can be rewritten as
\begin{align}
   p_{\tilde{\mathbf{y}}_{a,k},\mathbf{H}_{a,k}}\!(\tilde{\mathbf{x}}_{k}) = p_{\mathbf{H}_{a,k}}\!(\mathbf{H}_{a,k})p_{\tilde{\mathbf{y}}_{a,k}|\mathbf{H}_{a,k}}\!(\tilde{\mathbf{x}}_{k}),
\end{align}
where $p_{\tilde{\mathbf{y}}_{a,k}|\mathbf{H}_{a,k}}\!(\tilde{\mathbf{x}}_{k})$ is the conditional pdf of $\tilde{\mathbf{y}}_{a,k}$ and $\mathbf{H}_{a,k}$. As in~\cite{9882359}, to overcome the sensitivity of the Gaussian distribution to outliers, the heavy-tailed Student`s $t$ distribution is considered to describe $p_{\tilde{\mathbf{y}}_{a,k}|\mathbf{H}_{a,k}}\!(\tilde{\mathbf{x}}_{k})$, that is,
\begin{align}\label{eq:student}
    p_{\tilde{\mathbf{y}}_{a,k}|\mathbf{H}_{a,k}}\!(\tilde{\mathbf{x}}_{k})\!&\sim\!lst\left(\mathbf{H}_{a,k}\tilde{\mathbf{x}}_k, \sigma_{a}^2 \mathbf{I}, v\right),\nonumber\\
    &=\!\int_{0}^{\infty}\!\!\mathcal{N}\!\left(\!\mathbf{H}_{a,k}\tilde{\mathbf{x}}_{k}, \frac{\sigma_{a}^2 \mathbf{I}}{u}\!\right)\!\mathcal{G}\!\left(\!\frac{v}{2}, \frac{v}{2}\!\right)\!d u,
\end{align}
where $v$ is the degree of freedom, $u$ is an auxiliary hidden variable, and $\mathcal{G}$ is the gamma function. Thus, the latent variable is redefined as $\mathbf{h}_{a,k}=\left\{\mathbf{H}_{a,k}, u \right\}$, and the model parameters are given by $\mathbf{\vartheta}=\left\{\tilde{\mathbf{x}}_{k}, v\right\}$. Next, based on~\cite{art:em}, by introducing an arbitrary distribution $q(\mathbf{h}_{a,k})$ over the latent variable $\mathbf{h}_{a,k}$, $\mathcal{L}$ can be decomposed in two terms as
\begin{align}
	\mathcal{L}  \!=\!&\int q(\mathbf{h}_{a,k}) \log \left(p_{\tilde{\mathbf{y}}_{a,k},\mathbf{h}_{a,k}}\left(\mathbf{\vartheta}\right)\right) d \mathbf{h}_{a,k}, \nonumber\\
	 \!=\!&\int q(\mathbf{h}_{a,k}) \log\left( \frac{p_{\tilde{\mathbf{y}}_{a,k},\mathbf{h}_{a,k}}\left(\mathbf{\vartheta}\right)}{q(\mathbf{h}_{a,k})}\frac{q(\mathbf{h}_{a,k})}{p_{\mathbf{h}_{a,k}|\tilde{\mathbf{y}}_{a,k}}\left(\mathbf{\vartheta}\right)}\right) d \mathbf{h}_{a,k}, \nonumber\\
	\!=\!&\underbrace{\int q(\mathbf{h}_{a,k}) \log \left(\frac{p_{\tilde{\mathbf{y}}_{a,k},\mathbf{h}_{a,k}}\left(\mathbf{\vartheta}\right)}{q(\mathbf{h}_{a,k})}\right)d \mathbf{h}_{a,k}}_{\mathcal{F}_{\tilde{\mathbf{y}}_{a,k}}\!\left(q, \mathbf{\vartheta}\right)}+\mathrm{KLD},
 %\left(q(\mathbf{h}_{a,k}) \| p_{\mathbf{h}_{a,k}|\tilde{\mathbf{y}}_{a,k}}\left(\mathbf{\vartheta}\right)\right),
\end{align}

%& \geq \underbrace{\int q(\mathbf{h}_{a,k}) \log \frac{p\left(\mathbf{R}_{\mathrm{d}}, \mathbf{h}_{a,k} ; \mathbf{\vartheta}\right)}{q(\mathbf{h}_{a,k})} d \mathbf{h}_{a,k}}_{\mathcal{F}\left(\tilde{\mathbf{y}}_{a,k} ; q, \mathbf{\vartheta}\right)},

where $\mathrm{KLD}$ is given by
\begin{align}
	\mathrm{KLD}\!
 %\left(\!q(\mathbf{h}_{a,k}) \| p_{\mathbf{H}_{a,k}|\tilde{\mathbf{y}}_{a,k}}\!\left(\mathbf{\vartheta}\right)\!\right)\!\!=
 =\!\!\!\int\!\! q(\mathbf{h}_{a,k})\! \log\!\left(\! \frac{q(\mathbf{h}_{a,k})}{p_{\mathbf{h}_{a,k}|\tilde{\mathbf{y}}_{a,k}}\left(\mathbf{\vartheta}\right)}\!\right)\!d \mathbf{h}_{a,k},
\end{align}
which is the Kullback-Leiber divergence, and has the property of always being greater or equal to zero. Hence, the E-step and the M-step of the EM algorithm can be described only in terms of $\mathcal{F}_{\tilde{\mathbf{y}}_{a,k}}\!\left(q, \mathbf{\vartheta}\right)$ as presented next.
 \subsubsection{E-step} 
 $\mathcal{F}$ is maximized in terms of $q(\mathbf{h}_{a,k})$ for a fixed $\mathbf{\vartheta}$. Thus, after every $(a\!+\!1)$th iteration, $q(\mathbf{h}_{a,k})$ is updated as
\begin{align}
	q^{a+1}(\mathbf{h}_{a,k})=\arg \max _{q(\mathbf{h}_{a,k})} \mathcal{F}_{\tilde{\mathbf{y}}_{a,k}}\!\left(q, \mathbf{\vartheta}^a\right) .
\end{align}
 Given that $\mathbf{H}_{a,k}$ and $u$ are independent, $q$ can be factorized as $q(\mathbf{h}_{a,k})$ = $q(\mathbf{H}_{a,k})q(u)$. And based on~\cite{9882359}, $q(\mathbf{H}_{a,k})$ and $q(u)$ are optimal as
\begin{align}
	&q(\mathbf{H}_{a,k})\!\!\propto p_{\mathbf{H}_{a,k}}\!(\mathbf{H}_{a,k})\!\exp\!\left\{\!\mathbb{E}_{u}\!\left[\!\log\!\mathcal{N}\!\left(\!\mathbf{H}_{a,k}\tilde{\mathbf{x}}_{k}, \frac{\sigma_{a}^2 \mathbf{I}}{u}\!\right)\!\right]\!\right\},\label{eq:qhak} \\
	&q\!\left(u\right)\!\!\propto\mathcal{G}\!\!\left(\!\frac{v}{2}, \frac{v}{2}\right)\!\exp\!\left\{\!\mathbb{E}_{\mathbf{H}_{a,k}}\!\!\left[\!\log \mathcal{N}\!\left(\!\mathbf{H}_{a,k}\tilde{\mathbf{x}}_{k}, \!\frac{\sigma_{a}^2 \mathbf{I}}{u}\!\right)\!\right]\!\right\}\!.\label{eq:qu}
\end{align}
Accordingly, the distribution in \eqref{eq:qhak} have the mean and covariance given, respectively, by
\begin{align}
	\mathbb{E}_{\mathbf{H}_{a,k}}[q(\mathbf{H}_{a,k})] & \!=\!\!\left[\boldsymbol{\Omega}_{\mathbf{H}}\!\left(\frac{1}{\sigma_{\epsilon}^{2}} \Hat{\mathbf{H}}_{a,k}^{H}\!+\!\tilde{\mathbf{x}}_{k} \tilde{\mathbf{y}}_{a,k}^H \frac{\mathbb{E}_{u}\left[u\right]}{\sigma_{a}^2}\right)\!\right]^{H}\!,\label{eq:media}\\
	\boldsymbol{\Omega}_{\mathbf{H}_{a,k}} & =\left[\frac{1}{\sigma_{\epsilon}^{2}} \mathbf{I}+\tilde{\mathbf{x}}_{k}\tilde{\mathbf{x}}_{k}^{H} \frac{\mathbb{E}_{u}\left[u\right]}{\sigma_{a}^2}\right]^{-1}.\label{eq:rowcov}
	%\boldsymbol{\Sigma}_{\mathbf{H}}&=\mathbf{I}\label{eq:colcov}
\end{align}  
$q\left(u\right)$ on the other hand, can be approximated to $\mathcal{G}\!\left(\!\frac{v+2 N}{2},\!\frac{v+C}{2}\!\right)$, with
\begin{align}
	C\!=\!&\Bigg(\!\!\!\left(\tilde{\mathbf{y}}_{a,k}\!-\!\mathbb{E}_{\mathbf{H}_{a,k}}[q(\mathbf{H}_{a,k})] \tilde{\mathbf{x}}_{k}\right)^{H}\!\!\left(\tilde{\mathbf{y}}_{a,k}\!-\!\mathbb{E}_{\mathbf{H}_{a,k}}[q(\mathbf{H}_{a,k})] \tilde{\mathbf{x}}_{k}\right)\nonumber\\&+\tilde{\mathbf{x}}_{k}^{H}\left[\boldsymbol{\Omega}_{\mathbf{H}_{a,k}} M_{\mathrm{AP}}\right] \tilde{\mathbf{x}}_{k}\!\!\Bigg)\frac{1}{\sigma_{a}^2}.
\end{align}
\subsubsection{M-step}
For the variational $\mathrm{M}$-step, $\mathcal{F}$ is maximized in terms of $\mathbf{\vartheta}$ with fixed $q(\mathbf{h}_{a,k})$. Again, after every $(a\!+\!1)$th iteration, $\mathbf{\vartheta}^{a+1}$ is obtained by
\begin{align}
	\mathbf{\vartheta}^{a+1}=\arg \max _{\mathbf{\vartheta}} \mathcal{F}_{\tilde{\mathbf{y}}_{a,k}}\!\left(q^{a+1}, \mathbf{\vartheta}\right).
\end{align}

 Given \eqref{eq:media} and \eqref{eq:rowcov}, and after some mathematical manipulations,  $\mathcal{F}_{\tilde{\mathbf{y}}_{a,k}}\!\left(q^{a+1}, \mathbf{\vartheta}\right)$ can be written as $\sigma_{a}^2C.$
% \begin{align}\label{eq:newf}
%   \mathcal{F}_{\tilde{\mathbf{y}}_{a,k}}\!\left(q^{a+1}, \mathbf{\vartheta}\right)=\left(\tilde{\mathbf{y}}_{a,k}\!-\!\mathbb{E}_{\mathbf{H}_{a,k}}[q(\mathbf{H}_{a,k})] \tilde{\mathbf{x}}_{k}\right)^{H}\!\!\left(\tilde{\mathbf{y}}_{a,k}\!-\!\mathbb{E}_{\mathbf{H}_{a,k}}[q(\mathbf{H}_{a,k})] \tilde{\mathbf{x}}_{k}\right)\nonumber\\&+\tilde{\mathbf{x}}_{k}^{H}\left[\boldsymbol{\Omega}_{\mathbf{H}} M_{\mathrm{AP}}\right] \tilde{\mathbf{x}}_{k}
% \end{align}
Thus, the estimated transmitted interference signal, $\Hat{\mathbf{x}}_{k}$  can be recovered by minimizing $\mathcal{F}_{\tilde{\mathbf{y}}_{a,k}}\!\left(q^{a+1}, \mathbf{\vartheta}\right)$. 
However, to facilitate the analysis, a Cholesky factorization may be applied on $\boldsymbol{\Omega}_{\mathbf{H}_{a,k}}  M_{\mathrm{AP}}=\mathbf{V}_{\mathbf{H}}^{\top} \mathbf{V}_{\mathbf{H}}$, and $\Hat{\mathbf{x}}_{k}$ is obtained via the following optimization
\begin{align}\label{eq:optimalx}
	\Hat{\mathbf{x}}_{k} & \!=\!\min\!\left(\left\|\tilde{\mathbf{y}}_{a,k}\!-\!\mathbb{E}_{\mathbf{H}_{a,k}}[q(\mathbf{H}_{a,k})] \tilde{\mathbf{x}}_{k}\right\|^{2}\!\!+\!\left\|\mathbf{V}_{\mathbf{H}} \tilde{\mathbf{x}}_{k}\right\|^{2}\right).
\end{align}

The EM implementation is described in Algorithm 2.
\begin{algorithm}\label{alg:coalitional}
\caption{EM Algorithm}
\begin{algorithmic}[1]
\footnotesize		 \Initialize{$\mathbb{E}_{\mathbf{H}_{a,k}}^{(0)}[q(\mathbf{H}_{a,k})]\gets \hat{\mathbf{H}}_{a,k}$; $\Hat{\mathbf{x}}_{k}^{(0)}\gets \Hat{\mathbf{x}}_{k}^{\mathrm{ZF}}\ $; $a \gets 0$}

		\While{convergence = 0}
		%\State Calculate $q^{(a+1)}(u)$$=$$\mathcal{G}\!\left(\!\frac{v+2 N}{2},\!\left(\frac{v+C}{2}\right)^{a+1}\right)$ and $\mathbb{E}_u[u]$$=$$\left(\!\frac{v+2 N}{2}\!\right)\!/\!\left(\!\frac{v+C}{2}\!\right)^{(a+1)}$.
        \State Calculate $q^{(a+1)}(u)$ and $\mathbb{E}_u[u]^{(a+1)}$.
		\State Update $\mathbb{E}_{\mathbf{H}_{a,k}}^{(a+1)}$ according to \eqref{eq:media}, and $\boldsymbol{\Omega}^{(a+1)}_{\mathbf{H}}$ as in \eqref{eq:rowcov}.
		%\State Update $\mathbb{E}_{\mathbf{H}}^{(a+1)}$ according to \eqref{eq:media}.
		\State Do the Cholesky decomposition:
			$\boldsymbol{\Omega}_{\mathbf{H}}^{\mathrm{(a+1)}}\!M_{\mathrm{AP}}$$=$$\mathbf{V}_{\mathbf{H}}^{(a +1)T} \mathbf{V}_{\mathbf{H}}^{(a+1)}$.
	\State Solve the minimization problem in \eqref{eq:optimalx} to achieve $\Hat{\mathbf{x}}_k^{(a+1)}$ given $\mathbf{V}_{\mathbf{H}}^{(a+1)}$ and $\mathbb{E}_{\mathbf{H}_{a,k}}^{(a+1)}$.
	% \begin{align}
	% 	\hat{\mathbf{x}}^{(a+1)} & =\underset{\mathbf{x}}{\arg \min }\left\|\tilde{\mathbf{y}}_{a,k}-\mathbb{E}_{\mathbf{H}}^{(a+1)}[\mathbf{H}] \mathbf{x}\right\|^{2}+\left\|\mathbf{V}_{\mathbf{H}}^{(a+1)} \mathbf{x}\right\|^{2}.
	% \end{align}
		\If{$||\hat{\mathbf{x}}_k^{a}-\hat{\mathbf{x}}_k^{a-1}|| \leq cov$}
		\State convergence = 1
		%\EndIf
		\Else
		\State $a \gets a+1$;
		\State Return to Step 3.
		\EndIf
		\EndWhile
		\State \Return $\hat{\mathbf{x}}_k^{a}$.
\end{algorithmic}
\end{algorithm}

In Algorithm 2, $\hat{\mathbf{x}}_k^{\mathrm{ZF}}$ is computed via a zero-forcing (ZF) detector as 
\begin{align}
    \hat{\mathbf{x}}_k^{\mathrm{ZF}} = (\hat{\mathbf{H}}_{a,k}^H\hat{\mathbf{H}}_{a,k})^{-1}\hat{\mathbf{H}}_{a,k}^H\mathbf{y}_{a,k}.
\end{align}

Accordingly, after $\hat{\mathbf{x}}_k$ is obtained, its covariance matrix is computed as $\mathbf{R}_{\hat{\mathbf{x}}_k}$$=$$\hat{\mathbf{x}}_k\hat{\mathbf{x}}_k^H$, and the estimated beampattern for the $k$th transmitter AP is designed as 
\begin{align}
    \hat{\mathbf{B}}_k = \mathbf{a}^H(\theta_{z})\mathbf{R}_{\hat{\mathbf{x}}_k}\mathbf{a}(\theta_{z}), \forall k,
\end{align}
As illustrated in Fig.~\ref{fig:EM}, the adversary considers the maximum angle from $\hat{\mathbf{B}}_k$ as the estimated direction of the beam of AP $k$ to the target and defines a line of search inscribed in a square search area, having as many lines as number of transmitter APs. Next, the adversary divides the search area in cells of same size and verifies which of them are crossed by the computed direction lines. At the end, the cell crossed by the highest number of lines is set as the estimated position for the target. In case more than one cell contains the maximum number of crossing lines, the estimated position is selected randomly among the cells with maximum value.
\begin{figure}[t]
    \centering
    \includegraphics[scale=0.203]{Figures/Untitled Diagram-Page-2.drawio-4.pdf}
    \caption{Schematic of the process to compute the location of the target.}
    \label{fig:EM}
    \vspace{-0.52cm}
\end{figure}
Thus, the probability of detection for an observation $q$ is given by
\begin{align}
    P_{D,q} = \begin{cases}
      1 & \text{if $\hat{Pos_{t,q}}=Pos_t$}\\
      0 & \text{otherwise},
    \end{cases}       
\end{align}
where $Pos_t$ and $\hat{Pos}_t$ ar, respectively, the true and estimated position of the target. Hence the probability of detection for the total number of observations $Q$ is given by $P_D$$=$$\sum_{q=1}^{Q}P_{D,q}/Q$. 
\section{Numerical results and discussions}
For the numerical results, the search area is assumed to have 1000 m$\times$ 1000 m, with cells of 50 m$\times$ 50m and the coordinate $(0,0)$ positioned at the center of the search area. It is considered $N_{\mathrm{Tx}}$$=$$8$ transmitter APs serving $N_{\mathrm{Ue}}$$=$$4$ UEs with $M_{\mathrm{Ue}}$$=$$8$ antennas each. It is also assumed $N_{\mathrm{Rx}}$$=$$4$ sensing receiver APs, and both transmitter and receiver APS are equipped with $M_{\mathrm{AP}}$$=$$64$ antennas. The noise variance for each user, $\sigma_i^2$ , the communications SINR threshold, $\Gamma$, and the maximum transmit power per AP, $P_T$ are respectively set as -94 dBm, 3 dBm and 50 dBm. The carrier frequency is $\lambda_c$$=$$1.9$ GHz, and the variance of the RCS is $\sigma_{r,k}^2$$=$$10$ dBsm. In addition,  the average channel gains of the communications links are assumed to be determined by the pathloss, i.e., $\Omega_{i,k}$$=$$d_{i,k}^{-\varphi}$, where $d_{i,k}$ is the distance between the $k$th transmitter AP and UE $i$, and the path-loss exponent is set as $\varphi$$=$$3$. For algorithm 1,  the error tolerance $\epsilon$ is set as 0.1 and $p_{max}$$=$$10$ proved to render a sufficient number of iterations. Moreover, $\mathbf{W}^{(0)}$ and $\mathbf{u}_i(0)$ are chosen randomly. For algorithm 2, the variance of the channel estimation error is $\sigma_{\varepsilon}^2$$=$$10$ dBm and $cov$$=$$10^{-5}$. $P_D$ is computed over $Q$$=$$100$ channel realizations, and unless specified otherwise the coordinates of the target is $(-75,75)$, the receiver APs are at $(250, 250)$, $(-250,-250)$, $(-250,250)$ and $(250,-250)$. The UEs are positioned at $(300, 300)$, $(-300,-300)$, $(-300,300)$ and $(300,-300)$, , where the first one is considered as the adversary. Finally, the eight transmitter APS are at $(-500, -500)$, $(500,500)$, $(-500,500)$, $(500,-500)$, $(0,-500)$, $(0,500)$, $(-500,0)$ and $(500,0)$. 

\begin{figure}[t]
    \centering\includegraphics[scale=0.3]{Figures/PositionTarget.png}
    \caption{Probability of detection, $P_D$ vs. the position of the target for $M_{\mathrm{Ue}}$$=$$8, 16$ and $N_{\mathrm{Tx}}$$=$$4, 8$.}
    \label{fig:postarget}
       \vspace{-0.52cm}
\end{figure}

Fig.~\ref{fig:postarget} illustrates $P_D$ vs. the position of the target for $M_{\mathrm{Ue}}$$=$$8, 16$ and $N_{\mathrm{Tx}}$$=$$4, 8$. For this figure, the target's position is varied along the main diagonal of the search area, with the label of Fig.~\ref{fig:postarget} indicating the variation along the x-axis. Also, the coordinates of transmitter APS for the case with $N_{\mathrm{Tx}}$$=$$4$ are $(0,-500)$, $(0,500)$, $(-500,0)$ and $(500,0)$. Note that larger number of antennas at UEs increases the probability of detection by the adversary, as expected. Although the increase is noticeably more pronounced when the number of transmitter APs is also larger, as it increases the possibility that a bigger number of estimated direction lines converges to the cell that contains the target. Moreover, notice that the maximum $P_D$ for all cases occurs when the target is positioned closer to the center of the search area, which is also the closer point to the adversary location. In that case, the interference from the sensing signal at the adversary is higher, thus allowing for a better estimation from the EM algorithm. For the worst case, with $N_{\mathrm{Tx}}$$=$$8$ and $M_{\mathrm{Ue}}$$=$$16$, it leads to a $P_D$ higher than 0.7 which highlights the risk for the target's location privacy.    

\begin{figure}[t]
    \centering\includegraphics[scale=0.3]{Figures/Cellsize.png}
    \caption{Probability of detection, $P_D$ vs. the search area cell size for $M_{\mathrm{Ue}}$$=$$2,4, 8, 16$.}
    \label{fig:cellsize}
     \vspace{-0.52cm}
\end{figure}

In Fig.~\ref{fig:cellsize}, $P_D$ is evaluated for different sizes of cells and different values of the number of antennas at the UEs, $M_{\mathrm{Ue}}$$=$$2,4, 8, 16$. Validating the results observed in the previous figure, by increasing the number of antennas at the UEs, $P_D$ also increases. Note that for cases $M_{\mathrm{Ue}}$$=$$8$ and $M_{\mathrm{Ue}}$$=$$16$, there is a decrease on $P_D$ when the cell size is bigger than 100 m. It can be explained by the fact that there are fewer cells within the search area as the cell size increases, thus more cells might contain the maximum number of crossing lines, which reduces the probability that the chosen estimate position be the correct one. Even though, the $P_D$ remains relatively high, with the adversary identifying the correct cell in more than 50\% of the cases.

\begin{figure}[t]
    \centering\includegraphics[scale=0.3]{Figures/FigknownAPs.png}
    \caption{Probability of detection, $P_D$ vs. the number of known transmitter APs for $M_{\mathrm{Ue}}$$=$$2,4, 8, 16$.}
    \label{fig:knownaps}
     \vspace{-0.52cm}
\end{figure}

Fig.~\ref{fig:knownaps} illustrates $P_D$ vs. the number of known transmitter APs position from the adversary for $M_{\mathrm{Ue}}$$=$$2,4, 8, 16$. For this figure, we consider that the adversary receives data from all transmitter APs, but might not be capable to identify the location of all them. Note that, for $M_{\mathrm{Ue}}$$=$$8, 16$, with the knowledge of 6 transmitter APs, the adversary is already capable to achieve the maximum $P_D$ for both values of $M_{\mathrm{Ue}}$. Also, even with the position information of fewer APs, the adversary can still correctly detect the target in several opportunities. %This result portrays once more the serious risk of leakage of the sensing information in these scenarios.   
% \begin{table}[h]
% \centering
% \caption{Coordinates of the nodes for the simulation}
% \label{tab:parametros}
% \begin{tabular}{@{}lc@{}}
% \toprule
% %\multicolumn{2}{c}{\textbf{Room Specifications}}                                    \\ \midrule
% \multicolumn{1}{l|}{Transmitter APs}                         & (-500,-500), (500,500), (500,-500), (-500,500)\\ &(0, 500), (500, 0), (0,-500), (-500,0)               
% \\ 
% \multicolumn{1}{l|}{Receiver AP}                         & (250,250), (-250,-250), (-250,250), (250,-250)           \\ 
% \multicolumn{1}{l|}{UEs}                         & (300,300), (-300,-300), (-300,300), (300,-300)                   \\\bottomrule
% \end{tabular}
% \end{table}

\section{Conclusions}
In this paper, the capability of an UE, acting as adversary, to infer the position of a target in a multi-static sensing-centric ISAC system was investigated. For the adversary model, an EM algorithm was proposed to obtain an estimation of the transmitted signal, which was used to create replicas of the transmit beampattern of each transmitter AP. From the results, we showed that as more APs are deployed than transmitters, the adversary is capable to successfully identify the target location in a reasonable number of cases. Different from the mono-static scenario, where the adversary is capable to infer the angular position of the target, the multi-static sensing-centric scenario showed to be more critical in terms of privacy, since the adversary is capable of, with some accuracy, identify the exact location of the target. Such threats become even more critical as the network becomes denser.  
%\subsection{Future Works}


\section*{Acknowledgement}
This work has been supported by Academy of Finland, 6G Flagship program (Grant 346208) and FAITH Project (Grant 334280). Codes to reproduce results are available at: https://github.com/isabella-gomes/Globecom2023
\bibliographystyle{IEEEtran}
\bibliography{references.bib}
\end{document}